% AI-assisted working draft for team review, not an approved final section.
% Requires amsmath. Merge references.bib into the final report bibliography.
\section{Displacement, forces and matrix construction}
\label{sec:mechanics-matrix}

The preceding section introduced the nodal equilibrium system. We now
derive its entries from the displacement approximation and explain which
supports make the unknown displacements uniquely determined.

\subsection{Material response and force balance}
Our working baseline is a thin, homogeneous, isotropic sheet of constant
thickness $t>0$, loaded within its plane. We use plane stress, meaning
$\sigma_{zz}=\sigma_{xz}=\sigma_{yz}=0$, with free sheet faces. This
assumption permits thickness contraction; it does not prescribe zero
out-of-plane strain. We assume static loading, small displacement
gradients (including small rotations) and linear elastic behaviour.
The model excludes bending, plasticity and large rotations.

For displacement $\mathbf{u}=(u_x,u_y)^T$, stretching and changes of angle
give the dimensionless strain vector
\begin{equation}
 \varepsilon(\mathbf{u})=
 \begin{bmatrix}
 \partial u_x/\partial x\\
 \partial u_y/\partial y\\
 \partial u_x/\partial y+\partial u_y/\partial x
 \end{bmatrix},\qquad
 \sigma=D\varepsilon,\qquad
 D=\frac{E}{1-\nu^2}
 \begin{bmatrix}1&\nu&0\\\nu&1&0\\0&0&(1-\nu)/2\end{bmatrix}.
 \label{eq:mechanics-material}
\end{equation}
Here $\sigma=(\sigma_{xx},\sigma_{yy},\sigma_{xy})^T$, $E>0$ is Young's
modulus and $-1<\nu<1/2$ is Poisson's ratio for the underlying isotropic
material. The third strain component is engineering shear, twice the
tensor shear strain; this accounts for the shear entry of $D$~\cite{bleyer-plane-stress}.
Coordinates, displacements and thickness are in metres; stress and $E$
are in pascals. Numerical material values belong to each case specification.

To turn local force balance into nodal equations, multiply it by an
arbitrary virtual displacement $\mathbf{w}$ that vanishes in prescribed
components, integrate over the solid region $\Omega$ and integrate the
stress derivatives by parts. Internal and external work then satisfy
\begin{equation}
 \int_\Omega t\,\varepsilon(\mathbf{w})^TD\varepsilon(\mathbf{u})\,dA
 =\int_\Omega t\,\mathbf{w}^T\mathbf{g}\,dA
 +\int_{\Gamma_T}t\,\mathbf{w}^T\bar\tau\,ds.
 \label{eq:mechanics-work}
\end{equation}
Here $\mathbf{g}$ is body force per unit volume (N/m$^3$), $\Gamma_T$
is the boundary with prescribed traction, and $\bar\tau$ is that traction
(Pa), with horizontal and vertical components. It equals the stress tensor
times the outward unit normal. Unloaded boundary components have zero
traction~\cite{bleyer-plane-stress}.

\subsection{Triangle contributions and sparse assembly}
For general objects, the meshing specification uses ordered outer and hole
boundaries with constrained Delaunay triangulation through Python's
\texttt{triangle} interface to Triangle~\cite{shewchuk-triangle}.
The current rectangular benchmarks instead use a structured mesh split
into triangles. Both require positive-area triangles, matching shared
edges and empty holes. With $n$ nodes, the displacement vector is
$d=(u_{x,0},u_{y,0},\ldots,u_{x,n-1},u_{y,n-1})^T$.

On triangle $T_e$, let $(x_i,y_i)$, $i=1,2,3$, be its counterclockwise
vertices and $S_e>0$ its area. Linear shape functions interpolate
displacement as $\mathbf{u}_h=\sum_{i=1}^3N_i\mathbf{u}_i$,
where $N_i$ equals one at vertex $i$ and zero at the other two.
Writing $a_i=x_jy_k-x_ky_j$, $\beta_i=y_j-y_k$ and
$\eta_i=x_k-x_j$ for cyclic $(i,j,k)$ gives
$N_i=(a_i+\beta_i x+\eta_i y)/(2S_e)$. Therefore
\begin{equation}
 \varepsilon(\mathbf{u}_h)=B_ed_e,\qquad
 B_e=\frac{1}{2S_e}
 \begin{bmatrix}
 \beta_1&0&\beta_2&0&\beta_3&0\\
 0&\eta_1&0&\eta_2&0&\eta_3\\
 \eta_1&\beta_1&\eta_2&\beta_2&\eta_3&\beta_3
 \end{bmatrix},
 \label{eq:mechanics-B}
\end{equation}
where $d_e=(u_{x,1},u_{y,1},u_{x,2},u_{y,2},u_{x,3},u_{y,3})^T$
contains this triangle's nodal values. Thus strain is constant within each
triangle. Substituting the same interpolation for $\mathbf{u}$ and
$\mathbf{w}$ in Eq.~\eqref{eq:mechanics-work} gives
\begin{equation}
 K_e=\int_{T_e}tB_e^TDB_e\,dA=tS_eB_e^TDB_e,\qquad
 K=\sum_e P_e^TK_eP_e,
 \label{eq:mechanics-assembly}
\end{equation}
where $P_e$ selects the six entries $d_e=P_ed$ from $d$.
This is the constant-strain triangle construction~\cite{alberta-triangles}.
Each $6\times6$ contribution is added at its global node indices;
duplicate entries are summed and stored in compressed sparse row format.
Nodes without a common element do not couple directly, explaining sparsity.
The current implementation is \texttt{fem/elasticity.py}, function
\texttt{assemble}.\footnote{This is the team's implementation. The
course-provided assembly routine has not yet been located in this checkout;
permission to use the alternative remains unconfirmed.}

Loads follow from the right-hand side of Eq.~\eqref{eq:mechanics-work}.
For a constant boundary traction, write the line load $q=t\bar\tau$
(N/m). An edge of length $\ell$ contributes $\ell q/2$ to each endpoint's
two force components~\cite{alberta-triangles}. Body loads are integrated
with the same shape functions. The assembled vector $f$ has units N
and $K$ has units N/m; thickness is included once in each construction.
Refinement must preserve the physical loaded boundary and its traction.

\subsection{Supports, reduced system and reactions}
Let $C$ contain prescribed displacement components and $F$ its complement.
A node can be fixed in one direction while remaining free in the other.
Reordering $Kd=f+r$, with support forces zero on $F$, gives
\begin{equation}
 \begin{bmatrix}K_{FF}&K_{FC}\\K_{CF}&K_{CC}\end{bmatrix}
 \begin{bmatrix}d_F\\d_C\end{bmatrix}
 =\begin{bmatrix}f_F\\f_C+r_C\end{bmatrix},\qquad
 \underbrace{K_{FF}}_{A}\underbrace{d_F}_{x}
 =\underbrace{f_F-K_{FC}d_C}_{b}.
 \label{eq:mechanics-reduction}
\end{equation}
The first subscript of each block selects rows and the second selects
columns. Thus $A=K_{FF}$, $x=d_F$ and $b=f_F-K_{FC}d_C$ explicitly
identify the matrix, unknown and right-hand side of $Ax=b$. For fixed
supports $d_C=0$. After reconstructing $d$, the constrained rows recover
\begin{equation}
 r_C=K_{CF}d_F+K_{CC}d_C-f_C.
 \label{eq:mechanics-reactions}
\end{equation}
Keep the original $K$ and $f$, including loads applied at constrained
components. Their deletion would give incorrect reactions. Applied
forces plus reactions should balance both total force and moment;
$(Kd-f)_F$ measures the remaining free-equation imbalance.

\subsection{Matrix properties and sensitivity}
Equation~\eqref{eq:mechanics-assembly} makes $K$ symmetric. Its quadratic
form is
\begin{equation}
 d^TKd=\sum_e tS_e(B_ed_e)^TD(B_ed_e)\geq0,
 \label{eq:mechanics-energy}
\end{equation}
since $D$ is positive definite. Before supports, translations
$\mathbf{u}=(1,0)^T,(0,1)^T$ and infinitesimal rotation
$\mathbf{u}=(-y,x)^T$ have zero strain, so $K$ is singular.
For a conforming nondegenerate mesh on each connected solid component,
with no hinges joined only at a node, these are the only zero-energy
motions. Supports must remove all three on every component; merely
prescribing three arbitrary components need not suffice. Then any
nonzero free variation extended by zero on $C$ has positive energy,
so $A=K_{FF}$ is symmetric positive definite.

For this supported matrix, the 2-norm condition number is
\begin{equation}
 \kappa_2(A)=\frac{\lambda_{\max}(A)}{\lambda_{\min}(A)},\qquad
 \frac{\|\delta x\|_2}{\|x\|_2}
 \leq\kappa_2(A)\frac{\|\delta b\|_2}{\|b\|_2}.
 \label{eq:mechanics-conditioning}
\end{equation}
The eigenvalues $\lambda_{\min}$ and $\lambda_{\max}$ are positive;
the bound follows from $A\delta x=\delta b$ for fixed $A$ and nonzero
$b$. It bounds possible load sensitivity, rather than predicting every
perturbation's effect. We will use dense symmetric eigenvalues for small
systems and sparse extreme-eigenvalue estimates for larger ones
(\texttt{eigsh})~\cite{scipy-eigsh}, checking convergence and eigenpair
residuals and recording tolerances with the results. Geometry, mesh quality,
refinement and supports can affect conditioning. Small linear-system
residuals alone do not establish an accurate continuum approximation.
These properties provide the basis for the solver choices and comparison
protocol in the next section.
